\section{Conclusion}
\label{sec:conclusion}

We studied whether language models think better or worse when the user sounds skeptical or judgmental. Across 1,600 API calls on \gptfourone and \gptfouronemini, we found a weak, model-dependent sweet spot: mild skeptical questioning improves pooled no-hint accuracy for \gptfourone by 5.0 percentage points over neutral prompting, but harsher judgment does not add further benefit and the gain disappears under misleading-user pressure.

The main practical takeaway is that judgmental tone is not a reliable tool for better reasoning. Its most consistent effect is longer visible reasoning, not higher correctness, and the added length can coincide with worse answers. The misleading-hint results further show that the main failure mode is overcorrection and volatility rather than overt apology or direct agreement with the wrong user hint.

Future work should test this judgment-pressure ladder on additional model families, increase per-dataset sample sizes, and compare social judgment against non-social controls such as plain ``double-check carefully'' instructions. A richer annotation pass over changed answers would also help distinguish productive checking from unfaithful or unstable re-reasoning.
